% ===================================================================
%  3. TAULA — Haurtzaroa eta Gaztaroa.
%  Traduction 2026 d'apres texto/io, controlee sur texto/fr, texto/en
%  et texto/es. Consigne : voir texto/eu/10-taula-01.tex.
%
%  Deux scenes, et l'ido subdivise beaucoup plus que le francais.
%
%  LE BASQUE NOMME LE FRERE ET LA SŒUR SELON LE SEXE DE CELUI QUI
%  PARLE, et c'est le piege de ce tableau. « anaia » est le frere d'un
%  homme, « neba » le frere d'une femme ; « arreba » la sœur d'un
%  homme, « ahizpa » la sœur d'une femme. Quatre mots la ou toutes les
%  autres colonnes en ont deux. La sœur du nouveau-ne, qui est un
%  garcon, est donc « arreba » ; le frere de Berthe est « neba ».
%
%  LE FAC-SIMILE OUBLIE DE METTRE « julio » EN GRAS au bloc c2-01-1.
%  La source ne bouge pas ; gravuri/korekti.json recoit sa ligne
%  basque — « (uztaileko edo ».
% ===================================================================

%%K t03-apar-1 apar
\VUsaut{3.0mm}
\VUtitre{15.0pt}{60}{3. TAULA}
\VUsaut{1.6mm}
\VUtitre{11.4pt}{0}{Haurtzaroa eta Gaztaroa.}
\VUsaut{3.4mm}

%%K t03-apar-2 apar
(3. eta 4. taulak lau \VUgras{familia-eszenatan} \VUgras{eguraldiari},
\VUgras{adinari}, \VUgras{familiari}, \VUgras{jantziari} eta \VUgras{osasunari}
buruzko oinarrizko hitz-altxorra emateko prestatu dira.)
\VUsaut{4.0mm}

%%K t03-tit-1 sub
\VUcentre{12.0pt}{0}{\textit{Lehen agerraldia.}}
\VUsaut{3.0mm}

%%K t03-tit-2 sub
\VUpk{10.2pt}{40}{Bataioa herrian.}
\VUsaut{2.4mm}

%%K t03-tit-3 sub
\VUpk{10.2pt}{40}{Udaberria.}
\VUsaut{3.0mm}

%%K t03-c1-01-1 p
1. --- \VUgras{Udaberrian} gaude, \VUgras{martxoan}, \VUgras{apirilean} edo
\VUgras{maiatzean}, \VUgras{astearen} egunetako batean --- \VUgras{astelehenean},
\VUgras{asteartean} edo \VUgras{asteazkenean}. Goizeko \VUgras{hamarrak} dira.

%%K t03-c1-02-1 p
2. --- \VUgras{Eguraldia} bikaina da, \VUgras{airea} garbia. Eguzkiak distiratzen
du. \VUgras{Hodei} txiki batzuk \textsuperscript{(2)} igotzen dira \VUgras{zeru}
urdinean \textsuperscript{(1)}.

%%K t03-c1-03-1 p
3. --- \VUgras{Zuhaitzek} ia ez dute \VUgras{hostorik}. \VUgras{Makal}
lerdenak \textsuperscript{(3)} eta \VUgras{platano} garaia \textsuperscript{(17)}
kimuz bakarrik estalita daude oraindik.

%%K t03-c1-04-1 p
4. --- \VUgras{Enarak} \textsuperscript{(4)} itzuli dira eta oihu alaiez biraka
dabiltza \VUgras{orratzaren} \textsuperscript{(7)} inguruan:
\VUgras{kanpandorrearena} \textsuperscript{(5)} da, zeina herriko
\VUgras{elizaren} \textsuperscript{(6)} gainean altxatzen baita.

%%K t03-tit-4 sub
\VUsaut{3.4mm}
\VUpk{10.2pt}{40}{Eliza.}
\VUsaut{3.0mm}

%%K t03-c1-05-1 p
5. --- Oso xumea da eliza hau, \VUgras{atari} handi batekin \textsuperscript{(13)}
eta \VUgras{erloju} handi batekin \textsuperscript{(8)}. \VUgras{Orratz txikia}
\textsuperscript{(9)} X zenbakian dago: \VUgras{orduak} adierazten ditu.
\VUgras{Handia} \textsuperscript{(10)} XII.ean dago (eguerdia);
\VUgras{minutuak} adierazten ditu.

%%K t03-c1-06-1 p
6. --- Kanpandorrean \VUgras{kanpaiak} \textsuperscript{(11)} alai jotzen ari
dira. Teilatuaren gainean harrizko \VUgras{gurutze} txiki bat
\textsuperscript{(12)} dago.

%%K t03-c1-07-1 p
7. --- Elizak ez du atari handia bakarrik \textsuperscript{(13)}; alboko ate txiki
bat ere badu. Ate horretatik \VUgras{apaiza} \textsuperscript{(15)}, bere
\VUgras{sotanaz} jantzia, \VUgras{sakristiatik} \textsuperscript{(14)} irteten da
eta \VUgras{apaiz-etxera} \textsuperscript{(19)} abiatzen.

%%K t03-c1-08-1 p
8. --- Haren ondoan \VUgras{esnedun emakumea} \textsuperscript{(24)} dago, bere
\VUgras{astoarekin} \textsuperscript{(23)}. Hiritik itzultzen ari da, non
haurrentzat bere esne ona eraman baitu.

%%K t03-tit-5 sub
\VUsaut{4.0mm}
\VUpk{10.2pt}{40}{Fruta-baratzea.}
\VUsaut{3.4mm}

%%K t03-c1-09-1 p
9. --- Gris jauna (astoa) oso pozik dago goizeko paseo honekin. Gozamenez
arnasten du \VUgras{lorez} \textsuperscript{(20)} usaindutako airea, zeinak
alboko baratzeko \VUgras{fruta-arbolak} \textsuperscript{(18)} estaltzen
baititu. Arrosaz eta zuriz jantzita daude \VUgras{mertxikondo} horiek
\textsuperscript{(18)} eta \VUgras{abrikotondo} horiek \textsuperscript{(19)},
eta uzta ona agintzen dute.

%%K t03-c1-10-1 p
10. --- Bitartean \VUgras{erle} txikiak \textsuperscript{(22)},
\VUgras{erlauntzakoak} \textsuperscript{(21)}, hara doaz beren \VUgras{ezti}
gozoa biltzera, lanean burrunbaka.

%%K t03-c1-11-1 p
11. --- Baina zer gertatu da? Hara herriko haurrak korrika datoz eta
\VUgras{antzara} izutuak \textsuperscript{(25)} sakabanatzen dituzte. Segizio hau
elizatik ateratzen da, notarioaren semearen \VUgras{bataioaren} ondoren.

%%K t03-c1-12-1 p
12. --- Jakobe txikia bere \VUgras{sehaskan} \textsuperscript{(26)} esnatzen da
kanpaien soinu alaiaz, eta haren arreba Madalenak, zeinak segizioa ere ikusi nahi
baitu, amari eskatzen dio atarian orraz eta janz dezan.

%%K t03-c1-13-1 p
13. --- Amak baietz dio. Bere \VUgras{buruko zapia} \textsuperscript{(28)}, bere
\VUgras{jaka} \textsuperscript{(29)} eta bere \VUgras{mantala}
\textsuperscript{(30)} janzten ditu eta aulki baxu batean esertzen da atearen
aurrean. Madalenak \VUgras{gona} \textsuperscript{(31)} eta \VUgras{gerruntzea}
\textsuperscript{(32)} baino ez ditu soinean.

%%K t03-c1-14-1 p
14. --- Zein zoriontsu dagoen! Ahaztu egiten ditu bere lorerik maiteenak, bere
\VUgras{krabelinak} \textsuperscript{(34)}, zeinen gainean \VUgras{tximeletak}
\textsuperscript{(35)} pausatzen baitira, bere \VUgras{tulipak}
\textsuperscript{(36)}, bere \VUgras{ezkilaloreak} \textsuperscript{(37)}
\VUgras{loreontzietan} \textsuperscript{(38)} \VUgras{hormaren}
\textsuperscript{(33)} gainean, eta bere \VUgras{arrosak}
\textsuperscript{(39)}, zeinek hesi hain ederra osatzen baitute. Segizioko andreen
jantzi aberatsei begira geratzen da.

%%K t03-c1-15-1 p
15. --- Herriko mutilek gutxi begiratzen diete jantziei. Aurrera oldartzen dira
eta elkar bultzatzen dute \VUgras{gozokiak} \textsuperscript{(55)} edo
\VUgras{txanponak} \textsuperscript{(52)} lortzeko. Haietako bat negarrez ari da
\textsuperscript{(40)}.

%%K t03-c1-16-1 p
16. --- Beste batek \VUgras{txakur} baten \textsuperscript{(42)} hanka zapaldu du,
eta hura zaunkaka ihes doa eta \VUgras{oiloa} \textsuperscript{(43)} bere
\VUgras{txitekin} \textsuperscript{(44)} ihesi jartzen du.

%%K t03-tit-6 sub
\VUsaut{5.0mm}
\VUpk{10.2pt}{40}{Bataio-segizioa.}
\VUsaut{4.0mm}

%%K t03-c1-17-1 p
17. --- Segizioaren buruan \VUgras{inudea} \textsuperscript{(45)} doa. \VUgras{Txano}
eder bat \textsuperscript{(46)} darama, zinta luzeekin. \VUgras{Jaioberria}
\textsuperscript{(47)} darama, \VUgras{parpailaz} \textsuperscript{(48)} beterik.

%%K t03-c1-18-1 p
18. --- Inudearen atzetik \VUgras{aitabitxia} \textsuperscript{(49)} eta
\VUgras{amabitxia} \textsuperscript{(53)} datoz. Gozokiak eta txanponak banatzen
dituzte. Aitabitxiak \VUgras{eskularruak} \textsuperscript{(51)} eta
\VUgras{kapelu garaia} \textsuperscript{(50)} ditu. Amabitxiak \VUgras{lore-sorta}
eder bat \textsuperscript{(54)} darama eskuan.

%%K t03-c1-19-1 p
19. --- Haien atzetik haurraren \VUgras{osaba} \textsuperscript{(56)} eta
\VUgras{izeba} \textsuperscript{(58)} ikusten dira. Izebak \VUgras{kapela}
dotore bat \textsuperscript{(59)} darama; osabak \VUgras{lebita} beltza
\textsuperscript{(57)} eta txaleko zuria. Urrunago \VUgras{lehengusua}
\textsuperscript{(60)} eta \VUgras{lehengusina} \textsuperscript{(61)} datoz. Azken
honek eskutik darama jaioberriaren \VUgras{arreba} \textsuperscript{(62)}, Berta
txikia.

%%K t03-c1-20-1 p
20. --- Bertaren beste \VUgras{neba} \textsuperscript{(63)} atzetik doa.
\VUgras{Ahaide} bati \textsuperscript{(64)} ematen dio besoa. Haien atzetik
familiaren \VUgras{lagunak} \textsuperscript{(65)} datoz.

%%K t03-tit-7 sub
\VUsaut{4.6mm}
\VUpk{10.2pt}{40}{Begiluzeak.}
\VUsaut{3.4mm}

%%K t03-c1-21-1 p
21. --- Segizioaren ondoan \VUgras{eskale} bat \textsuperscript{(66)} dago, bere
\VUgras{makuluan} \textsuperscript{(67)} bermatua. Limosna eskatzen du.

%%K t03-c1-22-1 p
22. --- Beste aldean mutil oso \VUgras{adeitsu} bat \textsuperscript{(68)} dago.
Txapela altxatzen du eta \VUgras{«egun on»} esaten die ezagutzen dituenei.

%%K t03-c1-23-1 p
23. --- Herritarrak etxeetatik atera dira bataio-segizioa ikustera.
\VUgras{Nekazari} bat \textsuperscript{(69)} ikusten dugu, \VUgras{paperezko
txanoz} jantzia, eta haren ondoan \VUgras{nekazari emakume} bat
\textsuperscript{(70)}. Gero \VUgras{emakume zahar} bat \textsuperscript{(71)},
bere \VUgras{makilan} \textsuperscript{(72)} bermatua. Azkenik \VUgras{errotaria}
\textsuperscript{(73)}, \VUgras{feltrozko kapela} zabal batekin
\textsuperscript{(74)}, eta nekazari gazte bat \VUgras{eskalapoiekin}
\textsuperscript{(75)}, bere haurtxoari ibiltzen erakusten ari dena.

%%K t03-c1-24-1 p
24. --- Oso koadro bizia da.
\VUsaut{4.0mm}

%%K t03-tit-8 sub
\VUcentre{12.0pt}{0}{\textit{Bigarren agerraldia.}}
\VUsaut{3.0mm}

%%K t03-tit-9 sub
\VUpk{10.2pt}{40}{Herri-jaia.}
\VUsaut{2.4mm}

%%K t03-tit-10 sub
\VUpk{10.2pt}{40}{Ur-jokoak eta estropadak.}
\VUsaut{3.0mm}

%%K t03-c2-01-1 p
1. --- \VUgras{Uda} iritsi da. \VUgras{Ekaineko} (uztaileko edo \VUgras{abuztuko})
eguzki kiskalgarriak gazteak erakartzen ditu bainatzera eta txalupan ibiltzera.

%%K t03-c2-02-1 p
2. --- Ibaiaren eskuineko ertzean arraunlariak biluzten ari dira, ur-jokoetan eta
estropadetan parte hartzeko.

%%K t03-c2-03-1 p
3. --- Haietako batek bere \VUgras{oinetakoak} \textsuperscript{(1)} eranzten
ditu; lokarriak askatzen (botoiak askatzen) dizkie. Haren bi lagun prest daude
jada. \VUgras{Elastikoa} \textsuperscript{(2)} eta \VUgras{bainu-galtzak}
\textsuperscript{(3)} baino ez dituzte soinean orain. Berriketan ari dira,
lagunen zain. Feriaz eta beste ertzeko jaiaz ari dira.

%%K t03-c2-04-1 p
4. --- Lurrean, haien ondoan, lagunen \VUgras{arropa} datza: \VUgras{bota}
\VUgras{pare} bat \textsuperscript{(4)}, \VUgras{zapatilak}
\textsuperscript{(5)}, \VUgras{galtzerdiak} \textsuperscript{(6)},
\VUgras{feltrozko} \textsuperscript{(7)} eta \VUgras{lastozko}
\textsuperscript{(8)} kapelak eta \VUgras{gorbata} bat \textsuperscript{(9)}.

%%K t03-c2-05-1 p
5. --- Gazteetako batek txukun tolestu du bere \VUgras{jaka}
\textsuperscript{(10)}. Eskuoihal baten gainean jartzen du, zeinean haren
aldamenekoak berehala bilduko baititu beren bi \VUgras{jantziak}.

%%K t03-c2-06-1 p
6. --- Seigarren mutila berandutzen ari da. Buruan \VUgras{txapela}
\textsuperscript{(11)} du oraindik. Ez du erantzi ez \VUgras{alkandora}
\textsuperscript{(12)}, ez \VUgras{lepokoa} \textsuperscript{(13)}, ez
\VUgras{eskumuturrekoak} \textsuperscript{(14)}. Bere \VUgras{txalekoa}
\textsuperscript{(17)} eskuan darama eta oraindik \VUgras{galtzak}
\textsuperscript{(16)} eta \VUgras{tiranteak} \textsuperscript{(15)} ditu soinean.
Hurrengo lasterketara ez da iritsiko, seguru asko.

%%K t03-c2-07-1 p
7. --- Gazteen ondoan \VUgras{karduz} \textsuperscript{(18)} betetako bidezidor
bat urera jaisten da, non Jakobe zaharrak \VUgras{lezken} artean
\textsuperscript{(19)} \VUgras{su artifizialak} \textsuperscript{(20)} prestatzen
baititu, arratsean botatzeko.

%%K t03-tit-11 sub
\VUsaut{4.4mm}
\VUpk{10.2pt}{40}{Estropada.}
\VUsaut{3.4mm}

%%K t03-c2-08-1 p
8. --- Ibaian andre batzuk, aterkien azpian babestuta, \VUgras{txalupan}
\textsuperscript{(21)} dabiltza. Estropadari begira daude, eta oso jostagarria da.
\VUgras{Arraunontzi} bakoitzean \textsuperscript{(22)} lau \VUgras{arraunlarik}
\textsuperscript{(24)} indar guztiez arraun egiten dute, eta \VUgras{lemazainak}
\textsuperscript{(26)} \VUgras{lema} \textsuperscript{(25)} eusten dio eta
ontzi hauskorra gidatzen. \VUgras{Arraunek} \textsuperscript{(23)} neurriz jotzen
dute ura, eta txalupa arinak abiadura zorabiagarrian doaz.

%%K t03-c2-09-1 p
9. --- Gazte batzuk zubiaren pilarearen ondoan lehiatzen dira \VUgras{zutoin
sarituaren} \textsuperscript{(28)} sariagatik. Haietako bat kontuz aurreratzen da
haga malgu eta labainkorrean zehar, muturrean finkatutako \VUgras{banderatxoa}
\textsuperscript{(29)} hartzeko. Haren \VUgras{aurkari} zorigaiztokoa
\textsuperscript{(30)} urera erori da. Ertzera itzultzen ari da igerian.

%%K t03-c2-10-1 p
10. --- Mutil handienak \VUgras{ur-borrokan} \textsuperscript{(32)} ari dira
\VUgras{kaiaren} \textsuperscript{(33)} ondoan. \VUgras{Jendetza} handi bat
\textsuperscript{(31)} joko horiei guztiei begira dago, \VUgras{zubiaren}
\textsuperscript{(27)} \VUgras{barandetan} bermatuta eta \VUgras{ertzean} zehar
pilatuta. Begiluze batzuk, hala ere, buelta ematen dute \VUgras{zutoinaren}
\textsuperscript{(45)} jokoari begiratzeko edo \VUgras{udaletxeko segizioari}
begira geratzeko, zeina \VUgras{sari-banaketara} baitoa.

%%K t03-c2-11-1 p
11. --- Lehenik \VUgras{mutiko} batzuk \textsuperscript{(35)} korrika doaz
\VUgras{musikarien} \textsuperscript{(36)} aurretik; gero \VUgras{alkatea}
\textsuperscript{(37)}, haren laguntzaileak eta herriko \VUgras{udal-kontseilua}
datoz. Azkenak \VUgras{suhiltzaileak} \textsuperscript{(38)} doaz.

%%K t03-tit-12 sub
\VUsaut{5.0mm}
\VUpk{10.2pt}{40}{Feria.}
\VUsaut{3.6mm}

%%K t03-c2-12-1 p
12. --- \VUgras{Feria-plazan} \textsuperscript{(39)} jendetza handia dago. Jendea
\VUgras{barraka} desberdinak ikustera dator: \VUgras{zurezko zaldiak}
\textsuperscript{(40)}, \VUgras{zirkua} \textsuperscript{(41)}, non emanaldia hasi
baita jada; \VUgras{herri-dantzaldia} \textsuperscript{(42)}, non herriko gazteek
\VUgras{dantzatzeko} atsegina hartzen baitute.

%%K t03-c2-13-1 p
13. --- Barrenean \VUgras{zezen-plaza} ikusten da \textsuperscript{(44)}, non
\VUgras{matadore} espainiar ausartak \VUgras{zezen} amorratuarekin
\textsuperscript{(45)} borrokatzen baitu; gero \VUgras{errusiar mendiak}
\textsuperscript{(49)} eta \VUgras{tiro-lekua} \textsuperscript{(46)}, non
\VUgras{su-armak} erabiltzen eta \VUgras{ituan} tiro egiten trebatzen baitira.

%%K t03-c2-14-1 p
14. --- Ondoan \VUgras{barraka-antzokia} \textsuperscript{(47)} dago, non
\VUgras{komedia} eta \VUgras{drama} antzezten baitira. Oso hurbil
\VUgras{erraldoiaren} \textsuperscript{(48)} eta \VUgras{ipotxaren} barraka dago.
Lehenengoaren \VUgras{garaiera} bi metro eta hamabost da; bigarrena ozta-ozta da
haur bat bezalakoa.

%%K t03-c2-15-1 p
15. --- Azkenik, hona hemen \VUgras{animalia-erakusketa} handia
\textsuperscript{(50)}, bere \VUgras{basapiztiekin}, \VUgras{tigrearekin}
\textsuperscript{(51)} eta haren \VUgras{atzapar} indartsuekin,
\VUgras{lehoiarekin} \textsuperscript{(52)} eta haren \VUgras{adats} lodiarekin,
bi \VUgras{jirafarekin} \textsuperscript{(53)} eta \VUgras{elefantearekin}
\textsuperscript{(54)}, zeinak hain trebeki erabiltzen baitu bere
\VUgras{tronpa}.

%%K t03-c2-16-1 p
16. --- \VUgras{Hezitzailea} \textsuperscript{(55)}, animalia horiek hezten
dituena, barrakaren atean dago.
\VUsaut{6.0mm}
\VUfilet{22mm}
